% TOP_PROBLEM: {"id": 4800006, "title": "Multiple ergodic averages — Problem 6", "category_id": 11, "category": {"id": 11, "name": "dynamical_systems", "display_name": "Dynamical Systems", "description": "Problems about long-term behavior of deterministic systems, Hamiltonian dynamics, and periodic orbits.", "slug": "dynamical-systems", "order_index": 11, "created_at": "2026-07-31T15:26:25.671Z"}, "status": "partially_solved", "problem_number": "AMR-047-0006", "set_id": 13}
% FIRST_POSTED: 2026-10-01
% ENTRY_KIND: statement_only
% UnsolvedMath ID 4800006; source catalogue code AMR-047-0006
% !TeX program = lualatex
% Definition and sources only; this file is not an LLM solution attempt.
\documentclass[11pt,a4paper]{article}
\usepackage[margin=25mm]{geometry}
\usepackage{fontspec}
\setmainfont{lmroman10-regular.otf}[BoldFont=lmroman10-bold.otf,
 ItalicFont=lmroman10-italic.otf,BoldItalicFont=lmroman10-bolditalic.otf]
\usepackage{amsmath,amssymb,mathtools}
\usepackage{unicode-math}
\setmathfont{latinmodern-math.otf}
\AtBeginDocument{\let\setminus\smallsetminus}
\usepackage{xurl,hyperref}
\hypersetup{colorlinks=true,urlcolor=blue}
\setcounter{secnumdepth}{0}
\setlength{\parindent}{0pt}
\setlength{\parskip}{5pt}
\setlength{\emergencystretch}{3em}
\begin{document}
\section*{Multiple ergodic averages — Problem 6}\label{4800006}
{\small UnsolvedMath ID 4800006 · AMR-047-0006 · Dynamical Systems · AMR Open Problem Lists}
\subsection{Definitions and mathematical statement}
Let $a:\mathbb N\to\mathbb Z$ be an arbitrary sequence.
Say that it is good for two-fold convergence of powers
if for every probability space $(X,\mathcal B,\mu)$, every
invertible measure-preserving transformation $R$, every pair
of integers $r,s$, and every $f,g\in L^\infty(\mu)$, the averages
\[
     \frac1N\sum_{n=1}^N
          (f\circ R^{r a(n)})(g\circ R^{s a(n)})
\]
converge in the norm of $L^2(\mu)$ as $N\to\infty$.
Integer powers, including negative and zero powers, use
invertibility and the identity convention $R^0=\mathrm{id}$.

Prove that this hypothesis implies two-fold convergence for
commuting transformations: for every probability space,
every commuting invertible measure-preserving $T,S$, and
every bounded $f,g$, the averages
\[
     \frac1N\sum_{n=1}^N
          (f\circ T^{a(n)})(g\circ S^{a(n)})
\]
converge in $L^2(\mu)$.

No ergodicity, mixing, growth, monotonicity, or distinctness
assumption is made on the sequence or the transformations.
The limits may depend on the system and functions; no
particular limit formula is required. The hypothesis must
hold universally over all single transformations and all
integer pairs, rather than just for the pair of powers
$R,R^2$. Commuting convergence immediately implies the
hypothesis by choosing $T=R^r$, $S=R^s$; the question
asks whether this necessary condition is also sufficient.
\subsection{Short English statement}
Does universal mean convergence along powers of one transformation already imply mean convergence for any commuting pair?
\subsection{Sources}
\begin{itemize}
\item[S1] ULAM.AI and the UnsolvedMath Contributors, supplied problems.json, record 4800006 (AMR-047-0006), AMR Open Problem Lists. Catalogue identity and source transcription; CC BY 4.0.
\url{https://huggingface.co/datasets/ulamai/UnsolvedMath}.
\item[S2] Frantzikinakis - Some open problems on multiple ergodic averages (2016), Problem 6. Original source indexed by AMR-047-0006.
\url{https://arxiv.org/abs/1103.3808}.
\end{itemize}
\end{document}
